\section{Discussion and Conclusion}
\label{sec:discussion}

\subsection{Discussion}
\noindent\textbf{Task-level reliability.}
The two evaluated settings require later actions to depend on the outcome of earlier interaction with deformable material.
The results show that improving this task-level dependency is not achieved by changing the action generator alone.
ACT and DP frequently produce locally plausible behavior but fail to complete the full sequence, while the proposed representation supervision improves complete-task success across surgical color--layout conditions, daily object categories, and initial positions.
The daily failure-stage results further show that the largest difference occurs after the target has been exposed and separated, when the policy must use the resulting task state to complete transfer.

The factorial ablation clarifies the roles of the two representation components.
Task-query supervision provides the larger standalone gain, which supports directly preserving the relations needed for later decisions.
Semantic inputs alone provide a smaller improvement, showing that object-level visual structure does not by itself ensure that the policy uses the relevant task information.
Their combination performs best, indicating that structured observations and task-directed supervision serve complementary roles.

\noindent\textbf{Semantic supervision.}
The iterative annotation procedure expands a small reviewed set into a substantially larger supervision pool while maintaining high agreement on the reviewed masks.
The held-out flagged-frame rates improve from the initial to the final generation but are not monotonic across rounds.
This distinction is important because label coverage and label consistency are separate properties.
The current experiments evaluate the final annotation and policy pipeline, but do not isolate the contribution of quality screening from the effect of adding more semantic labels.

\noindent\textbf{Residual failures.}
All four daily-manipulation failures of Ours reach target separation but do not complete transfer into the goal region.
The current task descriptors encode visibility and image-plane geometry, but not contact mode, grasp stability, cloth tension, or force interaction.
These unmodeled factors provide concrete targets for extending the task representation.
Quality screening also detects temporal and structural anomalies rather than temporally stable segmentation errors.
The surgical simulation contains one timeout for Ours, while the baseline failures are dominated by workspace-limit and grasping failures.
These outcomes motivate explicit recovery supervision and representations that include contact state.

\begin{figure}[t]
\centering
\fbox{\parbox[c][2.8cm][c]{0.92\columnwidth}{\centering\small
(a) Persistent perception error: RGB, semantic prediction, and reviewed mask\\[4pt]
(b) Correct task descriptors but incomplete transfer: relation state and action trace\\[4pt]
(c) Surgical timeout: exposure, grasp, and transfer stages}}
\caption{Representative residual failure cases to be completed with synchronized observations, task descriptors, and action traces.}
\label{fig:failurecases}
\end{figure}

\noindent\textbf{Scope of the evidence.}
The evaluation uses fixed expert demonstrations and one trained checkpoint per configuration.
The crayon tests measure transfer to an object category absent from the policy demonstrations, but do not establish zero-shot transfer across tasks or robot platforms.
The surgical experiments use tissue-mimicking material and do not constitute clinical validation.
Ours also uses offline semantic supervision beyond the action demonstrations available to ACT and DP.
The theoretical analysis identifies conditions connecting action error, representation coverage, and task failure, while the experiments test the corresponding mechanisms rather than the numerical tightness of the bounds.

\subsection{Conclusion}
This work introduced a task-aware imitation-learning framework for bimanual sequential manipulation from fixed demonstrations.
A quality-aware annotation procedure expands semantic supervision, and supervised task queries preserve target visibility and interaction geometry in the action-generating representation.
The complete method succeeds in 19 of 20 surgical-simulation trials and 92 of 96 daily-manipulation trials, compared with 8 of 20 and 54 of 96 for ACT and 11 of 20 and 50 of 96 for Diffusion Policy.
In the daily factorial study, task-query supervision alone reaches 72 of 96 trials, semantic inputs alone reach 60 of 96, and their combination performs best.
These results demonstrate that explicitly preserving task-relevant state can improve complete-task reliability without additional policy interaction data or an explicit skill selector.
The remaining failures point to contact-aware task descriptors and recovery-focused supervision as the next steps.
